\begin{name}
{\tenchude}{\tendethi}{LỚP TOÁN THẦY PHÁT}{\thoigian}
\end{name}
\Opensolutionfile{ans}[ans/ansBTTeX1]
\begin{ex}%[2D3Y1-1]%
Họ tất cả các nguyên hàm của hàm số  $ f(x)=4^x $ là
\choice
{\True $ \dfrac{4^x}{\ln 4} $}
{$ 4^{x+1}+C $}
{$ \dfrac{4^{x+1}}{x+1} $}
{$ 4^x\ln 4+C $}
\loigiai{
Ta có $ \displaystyle\int 4^x\mathrm{\,d}x=\dfrac{4^x}{\ln 4}$.
}
\end{ex}

\begin{ex}%[2D3Y1-1]%
Cho hàm số $ f(x)=x^2+2$. Khẳng định nào dưới đây đúng?
\choice
{$\displaystyle\int{f(x)\,\mathrm{d}x=2x+C}$}
{\True $\displaystyle\int{f( x )\,\mathrm{d}x=\dfrac{x^3}{3}+2x+C}$}
{$\displaystyle\int{f( x )\,\mathrm{d}x=x^2+2x+C}$}
{$\displaystyle\int{f( x )\,\mathrm{d}x=x^3+2x+C}$}
\loigiai{
Ta có $\displaystyle\int f(x)\,\mathrm{d}x= \displaystyle\int\left( x^2+2 \right)\,\mathrm{d}x=\dfrac{x^3}{3}+2x+C$.
}
\end{ex}

\begin{ex}%[Vanle Vo]%[2D3B1-1]
Tìm họ nguyên hàm của hàm số $f(x)=\dfrac{1}{2x-1}$.
\choice
{$\displaystyle\int f(x){~\rm d}x=\ln \left|2x-1\right|+C$}
{\True $\displaystyle\int f(x){~\rm d}x=-\ln \left|2x-1\right|+C$}
{$\displaystyle\int f(x){~\rm d}x=\dfrac{1}{2}\ln \left|2x-1\right|+C$}
{$\displaystyle\int f(x){~\rm d}x=2\ln \left|2x-1\right|+C$}
\loigiai{
$\displaystyle\int f(x){~\rm d}x=\displaystyle\int\dfrac{1}{2x-1}{~\rm d}x=\dfrac{1}{2}\ln \left|2x-1\right|+C$
}
\end{ex}

\begin{ex}%[2D3B1-1]
Họ nguyên hàm của hàm số $f(x)=4x^3-2020$ là
\choice
{\True $x^4-2020x+C$}
{$12x^3+C$}
{$x^4+C$}
{$4x^3-2020x+C$}
\loigiai{
Ta có $\displaystyle\int (4x^3-2020)\mathrm{\,d}x=4 \cdot \dfrac{x^4}{4}-2020x+C=x^4-2020x+C$.
}
\end{ex}

\begin{ex}%[2D3B1-2]%
Cho $\int x\sqrt{x} \mathrm{d}x$ bằng
\choice
{$\dfrac{1}{3}x^3 \sqrt{x}+C$}
{$\dfrac{1}{2}x^2\sqrt{x}+C$}
{$\dfrac{5}{2}x^2\sqrt{x}+C$}
{\True $\dfrac{2}{5}x^2\sqrt{x}+C$}
\loigiai {
Đặt $\sqrt{x}=t \Rightarrow t^2=x \Rightarrow 2t\mathrm{d}t=\mathrm{d}x$.\\
Ta có $\int x\sqrt{x} \mathrm{d}x=\int t^2\cdot t\cdot 2t\mathrm{d}t=2\int t^4\mathrm{d}t=\dfrac{2}{5}t^5+C=\dfrac{2}{5}\left( \sqrt{x}\right) ^5+C=\dfrac{2}{5}x^2\sqrt{x}+C$.
}
\end{ex}

\begin{ex}%[2D3B1-3]%
Họ nguyên hàm của hàm số$f(x)=2x( \sin x+1 )$là
\choice
{$x^2+2x\cos x-2\sin x+C$}
{$x^2-2x\cos x-2\sin x+C$}
{$x^2( x-\cos x )+C$}
{\True $x^2-2x\cos x+2\sin x+C$}
\loigiai{
Ta có $\displaystyle \int{f( x )\mathrm{\,d}x}=\displaystyle \int{2x( \sin x+1 )\mathrm{\,d}x}$.\\
Đặt $\heva{&u=2x \\&\mathrm{\,d}v=( \sin x+1 )\mathrm{\,d}x}$ $\Rightarrow \heva{&\mathrm{\,d}u=2\mathrm{\,d}x \\&v=-\cos x+x.}$\\
Suy ra
\begin{eqnarray*}
\displaystyle \int{f(x)\mathrm{\,d}x} &=& \displaystyle \int{2x(\sin x+1)\mathrm{\,d}x}\\
&=& 2x( -\cos x+x )-2\displaystyle \int{(-\cos x+x)\mathrm{\,d}x}\\
&=& -2x\cos x+2x^2+2\sin x-x^2+C\\
&=& x^2-2x\cos x+2\sin x+C.
\end{eqnarray*}
}
\end{ex}

\begin{ex}%[2D3Y2-1]%
Cho $a<b<c,\displaystyle\int\limits_a^b f(x)\mathrm{\,d}x=5$ và $\displaystyle\int\limits_c^b f(x)\mathrm{\,d}x=2$. Tính $\displaystyle\int\limits_a^c f(x)\mathrm{\,d}x$.
\choice
{\True $\displaystyle\int\limits_a^c f(x)\mathrm{\,d}x=3$}
{$\displaystyle\int\limits_a^c f(x)\mathrm{\,d}x=-2$}
{$\displaystyle\int\limits_a^c f(x)\mathrm{\,d}x=1$}
{$\displaystyle\int\limits_a^c f(x)\mathrm{\,d}x=7$}
\loigiai{
Ta có $\displaystyle\int\limits_a^b f(x)\mathrm{\,d}x=\displaystyle\int\limits_a^c f(x)\mathrm{\,d}x+\displaystyle\int\limits_c^b f(x)\mathrm{\,d}x\Rightarrow 5=\displaystyle\int\limits_a^c f(x)\mathrm{\,d}x+2\Rightarrow\displaystyle\int\limits_a^c f(x)\mathrm{\,d}x=3$.\\
Vậy $\displaystyle\int\limits_a^c f(x)\mathrm{\,d}x=3$.}
\end{ex}

\begin{ex}%[2D3Y2-1]
Nếu $\displaystyle\int\limits_0^3 f(x)\mathrm{\,d}x=3,\displaystyle\int\limits_3^5 f(x)\mathrm{\,d}x=7$ thì $\displaystyle\int\limits_0^5 f(x)\mathrm{\,d}x$ bằng
\choice
{$7$}
{$4$}
{\True $10$}
{$-4$}
\loigiai{
Ta có $\displaystyle\int\limits_0^5 f(x)\mathrm{\,d}x=\displaystyle\int\limits_0^3 f(x)\mathrm{\,d}x+\displaystyle\int\limits_3^5 f(x)\mathrm{\,d}x=3+7=10$.
}
\end{ex}

\begin{ex}%[2D3B2-1]%
Gọi ${{z}_{1}}$ và ${{z}_{2}}$ là hai nghiệm phức của phương trình ${{z}^{2}}+z+3=0$. Khi đó $\left| {{z}_{1}} \right|+\left| {{z}_{2}} \right|$ bằng
\choice
{$3$}
{$6$}
{$\sqrt{3}$}
{\True $2\sqrt{3}$}
\loigiai{
Ta có ${{z}^{2}}+z+3=0\Leftrightarrow {{\left( z+\frac{1}{2} \right)}^{2}}=-\frac{11}{4}\Leftrightarrow {{\left( z+\frac{1}{2} \right)}^{2}}={{\left( \frac{\sqrt{11}i}{2} \right)}^{2}}\Leftrightarrow z=-\frac{1}{2}\pm \frac{\sqrt{11}}{2}i$.
Do đó $\left| {{z}_{1}} \right|+\left| {{z}_{2}} \right|=2\sqrt{{{\left( \frac{1}{2} \right)}^{2}}+{{\left( \frac{\sqrt{11}}{2} \right)}^{2}}}=2\sqrt{3}$
}
\end{ex}

\begin{ex}%[2D3B2-2]% ::Cau 34::

Cho $\int\limits_1^{13}{f\left( x \right)\mathrm{\,d}x}=2019$. Tính $\int\limits_0^4{f\left( 3x+1 \right)\mathrm{\,d}x}$.
\choice
{\True $673$}
{ $-2019$}
{ $2019$}
{ $6057$}
\loigiai{
Gọi $I=\int\limits_0^4{f\left( 3x+1 \right)\mathrm{\,d}x}$.\\
Đặt $t=3x+1\Rightarrow \mathrm{\,d}t=3\mathrm{\,d}x$.\\
Đổi cận
$\left[ \begin{aligned}
x=0\Rightarrow&\ t=1 \\
x=4\Rightarrow&\ t=13
\end{aligned}\right. .$\\
Suy ra $I=\int\limits_1^{13}{f\left( t \right)\text{.}\dfrac{1}{3}\mathrm{\,d}t}=\dfrac{1}{3}.2019=673.$}
\end{ex}

\begin{ex}%[2D3B2-2]%
Xét tích phân $I=\displaystyle\int\limits_0^4\mathrm{e}^{\sqrt{2x+1}}\mathrm{d}x$, nếu đặt $u=\sqrt{2x+1}$ thì $I$ bằng
\choice
{$\dfrac{1}{2}\displaystyle\int\limits_1^3u\mathrm{e}^u\mathrm{d}u$}
{$\displaystyle\int\limits_0^4u\mathrm{e}^u\mathrm{d}u$}
{\True $\displaystyle\int\limits_1^3u\mathrm{e}^u\mathrm{d}u$}
{$\dfrac{1}{2}\displaystyle\int\limits_1^3\mathrm{e}^u\mathrm{d}u$}
\loigiai{
Đặt $u=\sqrt{2x+1}\Rightarrow x=\dfrac{u^2-1}{2}\Rightarrow\mathrm{d}x=u\mathrm{d}u$.\\
Đổi cận: $x=0\Rightarrow u=1$, $x=4\Rightarrow u=3$.\\
Do đó $I=\displaystyle\int\limits_1^3 u \mathrm{e}^u \mathrm{d}u$.
}
\end{ex}

\begin{ex}%[2D3B2-3].
Nếu $\displaystyle\int_1^a\ln x \mathrm{\, d}x=1+2 a$ với $a>1$ thì $a$ thuộc khoảng nào sau đây?
\choice
{\True $(18; 21)$}
{$(1; 4)$}
{$(11; 14)$}
{$(6; 9)$}
\loigiai{
Đặt $\heva{& u=\ln x\Rightarrow \mathrm{\,d}u=\dfrac{1}{x}\mathrm{\,d}x \\ & \mathrm{\,d}v= \mathrm{\,d}x\Rightarrow v=x.}$\\
Khi đó $\displaystyle\int_1^a\ln x \mathrm{\,d}x=x\cdot\ln x\bigg|_1^a-\displaystyle\int_1^a x\cdot\dfrac{1}{x} \mathrm{\,d}x=a\ln a-x\bigg|_1^a=a\ln a-a+1$.\\
Theo giả thiết, ta có $a\ln a-a+1=1+2a\Leftrightarrow \ln a=3\Leftrightarrow a=\mathrm{e}^3\in(18;21)$.
}
\end{ex}

\begin{ex}%[2D3B2-3].
Tích phân $I=\displaystyle\int_1^2 x \mathrm{e}^x \mathrm{\,d}x$ bằng
\choice
{\True $\mathrm{e}^2$}
{$- \mathrm{e}^2$}
{$\mathrm{e}$}
{$3 \mathrm{e}^2-2 e$}
\loigiai{
Đặt $\heva{& u=x\Rightarrow \mathrm{\,d}u=\mathrm{\,d}x \\ & \mathrm{\,d}v=\mathrm{e}^x \mathrm{\,d}x\Rightarrow v=\mathrm{e}^x.}$\\
Khi đó $I=x\mathrm{e}^x\bigg|_1^2-\displaystyle\int_1^2 \mathrm{e}^x \mathrm{\,d}x=2\mathrm{e}^2-\mathrm{e}-\mathrm{e}^x\bigg|_1^2=2\mathrm{e}^2-(\mathrm{e}^2-2\mathrm{e})=\mathrm{e}^2$.
}
\end{ex}

\begin{ex}%[Dự án tex hệ thống các dạng toán thường gặp LVD]%[Quách Tấn Phát]%[2D3Y3-3]%
Cho hình $\left(D\right)$ giới hạn bởi các đường $y=f(x)$, $y=0$, $x=\pi$, $x=\mathrm{e}$. Quay $\left(D\right)$ quanh trục $Ox$ ta được khối tròn xoay có thể tích $V$. Khi đó $V$ được xác định bằng công thức
\choice
{$V=\pi \displaystyle\int\limits_{\mathrm{e}}^{\pi}\Big|f(x)\Big|\mathrm{\,d}x$}
{$V=\pi \displaystyle\int\limits_{\pi}^{\mathrm{e}}f^{2}(x)\mathrm{\,d}x$}
{\True $V=\pi \displaystyle\int\limits_{\mathrm{e}}^{\pi}f^{2}(x)\mathrm{\,d}x$}
{$V=\displaystyle\int\limits_{\mathrm{e}}^{\pi}\Big|f(x)\Big|\mathrm{\,d}x$}
\loigiai{
Ta có $\mathrm{e} < \pi$, ta áp dụng công thức, thể tích cần tính bằng
$$V=\pi \displaystyle\int\limits_{\mathrm{e}}^{\pi}f^{2}(x)\mathrm{\,d}x\text{.}$$ }
\end{ex}

\begin{ex}%[2D3Y3-1].
\immini
{
Cho đồ thị hàm số $y=f(x)$. Diện tích hình phẳng $S$ (phần tô đậm trong hình) bằng
\choice
{\True $\displaystyle\int_{-2}^0 f(x) \mathrm{\,d}x-\displaystyle\int_0^2 f(x) \mathrm{\,d}x$}
{$\displaystyle\int_{-2}^0 f(x) \mathrm{\,d}x+\displaystyle\int_0^2 f(x) \mathrm{\,d}x$}
{$\left|\displaystyle\int_{-2}^0 f(x) \mathrm{\,d}x+\displaystyle\int_0^2 f(x) \mathrm{\,d}x\right|$}
{$\displaystyle\int_{-2}^2 f(x) \mathrm{\,d}x$}
}
{
\begin{tikzpicture}[scale=1,font=\footnotesize,line join = round, line cap = round,>=stealth]
\def\hsf(#1){0.5*(#1)^3-2*(#1)}
\draw[->] (-2.5,0)--(0,0)node[below left]{$O$}--(3,0)node[below]{$x$};
\draw[->] (0,-2)--(0,2.5)node[left]{$y$};
\draw[samples=100,domain=-2.38:2.38,smooth] plot (\x, {\hsf(\x)});
\draw[pattern = north east lines,opacity=0.5,draw=none] plot[domain=-2:2] (\x, {\hsf(\x)})--cycle;
\path (-2,0)node[below left]{$-2$} (2,0)node[below right]{$2$} (-2.3,{\hsf(-2.3)})node[right]{$y=f(x)$};
\end{tikzpicture}
}
\loigiai{
Từ hình vẽ ta có diện tích bằng $\displaystyle\int_{-2}^0 f(x) \mathrm{\,d}x-\displaystyle\int_0^2 f(x) \mathrm{\,d}x$.
}
\end{ex}

\begin{ex}%[2D3Y3-1]%
Cho hàm số $y=f\left( x \right)$ xác định và liện tục trên đoạn $\left[ a;b \right]$. Diện tích hình phẳng giới hạn bởi đồ thị hàm số $y=f\left( x \right)$, trục hoành và hai đường thẳng $x=a$; $x=b$ được tính theo công thức
\choice
{ $S=\left| \int\limits_a^b{f\left( x \right)\mathrm{\, d}x} \right|$}
{ $S=\int\limits_a^b{f\left( x \right)\mathrm{\, d}x}$}
{ $S=\int\limits_b^a{\left| f\left( x \right) \right|\mathrm{\, d}x}$}
{\True $S=\int\limits_a^b{\left| f\left( x \right) \right|\mathrm{\, d}x}$}
\loigiai{
Công thức tổng quát tính diện tích hình phẳng là $S=\int\limits_a^b{\left| f\left( x \right) \right|\mathrm{\, d}x}$.}
\end{ex}

\begin{ex}%[Câu 17]%[2D3B3-1]%
Phần hình phẳng $(H)$ được gạch chéo trong hình vẽ dưới đây được giới hạn bởi đồ thị hàm số $y=f(x)$, $y=x^2+4x$ và hai đường thẳng $x=-2$; $x=0$.
\begin{center}
\begin{tikzpicture}[>=stealth,x=1cm,y=1cm,scale=0.4]
\draw[->] (-5,0) -- (3,0) node[below] {\scriptsize $x$};
\draw[->] (0,-5) -- (0,5) node[left] {\scriptsize $y$};
\fill[pattern=north east lines]plot[domain=-2:0](\x,{(\x)^2+4*\x})--(0,2)--plot[domain=0:-2](\x,{-(\x)^2+2})--cycle;
\draw[dashed] (0,0)node[below right]{\scriptsize $O$} (-2,0)node[above]{\scriptsize $-2$}--(-2,-4) (.75,4.5)node[right]{\scriptsize $y=x^2+4x$} (2.6,-4.5)node[right]{\scriptsize $y=f(x)$};
\clip (-5,-5)rectangle(3,5);
\draw[thick,samples=150,smooth,domain=-5:5] plot(\x,{(\x)^2+4*\x});
\draw[thick,samples=150,smooth,domain=-5:5] plot(\x,{-(\x)^2+2});
\end{tikzpicture}
\end{center}
Biết $\displaystyle \int\limits_{-2}^0 f(x)\mathrm{\,d}x=\dfrac{4}{3}$. Diện tích hình $(H)$ là
\choice
{$\dfrac{7}{3}$}
{$\dfrac{16}{3}$}
{$\dfrac{4}{3}$}
{\True $\dfrac{20}{3}$}
\loigiai{
Diện tích hình $(H)$ là
\begin{eqnarray*}
S&=&\displaystyle\int\limits_{-2}^0\left[f(x)-(x^2+4x)\right]\mathrm{\,d}x
\\&=&\displaystyle\int\limits_{-2}^0 f(x)\mathrm{\,d}x-\int\limits_{-2}^0(x^2+4x)\mathrm{\,d}x\\&=&\dfrac{4}{3}-\left(\dfrac{x^3}{3}+2x^2\right)\left| \begin{aligned}
&^0 \\
& _{-2} \\
\end{aligned}\right.\\&=&\dfrac{4}{3}+\dfrac{(-2)^3}{3}+2(-2)^2=\dfrac{20}{3}.
\end{eqnarray*}
Vậy diện tích hình $(H)$ là $S=\dfrac{20}{3}$.}
\end{ex}

\begin{ex}%[2D3B3-3]%[Võ Thị Thùy Trang, LaTeX hóa Đề cương HK2 - Lớp 12 - Lê Văn Đoàn]%17
\immini{
Thể tích của vật thể tròn xoay thu được khi quay hình phẳng (phần gạch sọc của hình vẽ) xung quanh trục hoành bằng
\choice
{$\dfrac{12\pi}{5}$}
{$\dfrac{53\pi}{15}$}
{\True $\dfrac{153\pi}{5}$}
{$\dfrac{31\pi}{13}$}
}
{\begin{tikzpicture}[scale=.7,>=stealth, font=\footnotesize, line join=round, line cap=round]
\def\a{-1} \def\b{0} \def\c{5} % Hệ số
\def\xmin{-3} \def\xmax{5}
\def\ymin{-1} \def\ymax{7}
%   \draw[color=gray!50,dashed] (\xmin,\ymin) grid (\xmax,\ymax);
\draw[->] (\xmin,0)--(\xmax,0) node [below]{$x$};
\draw[->] (0,\ymin)--(0,\ymax) node [left]{$y$};
\node at (0,0) [below left]{$O$};
\node at (-2,5.5) [right]{$y=3-x$};
\node at (1,4.55) [right]{$y=5-x^2$};
\draw[dashed] (-1,0) circle (1pt)--(-1,4) circle (1pt);
\draw[dashed] (2,0) circle (1pt)--(2,1) circle (1pt);
\clip (\xmin+0.1,\ymin+0.1) rectangle (\xmax-0.5,\ymax-0.1);
\draw[smooth,samples=300] plot(\x,{\a*(\x)^2+\b*(\x)+\c});
\draw[smooth,samples=300] plot(\x,{(-1)*(\x)+3});
\fill[pattern=north east lines,opacity=0.8] plot[domain=-1:2](\x,{\a*(\x)^2+\b*(\x)+\c})--plot(\x,{(-1)*(\x)+3})--cycle;
\end{tikzpicture}}
\loigiai{
Phương trình hoành độ giao điểm $3-x=5-x^2 \Leftrightarrow x^2-x-2=0 \Leftrightarrow \hoac{&x=-1\\&x=2.}$\\
Thể tích của vật thể tròn xoay cần tìm là\\ $V=\pi\left|\displaystyle\int\limits_{-1}^2 (5-x^2)^2 \mathrm{\,d}x-\displaystyle\int\limits_{-1}^2 (3-x)^2 \mathrm{\,d}x\right|=\dfrac{153\pi}{5}$.
}
\end{ex}

\begin{ex}%[2D4Y1-1]%
Cho số phức $z=4-3 i$ mô-đun của nó tương ứng là
\choice
{$|z|=1$}
{$|z|=25$}
{\True$|z|=5$}
{$|z|=\sqrt{5}$}
\loigiai{
Ta có $z=\sqrt{4^2+(-3)^2}=5$.
}
\end{ex}

\begin{ex}%[2D4Y1-1]%
Cho số phức $z$ và $w$ có điểm biểu diễn trong mặt phẳng $Oxy$ lần lượt là $M\left( 2\,;\,1 \right)$ và $N\left( 1\,;\,2 \right)$. Tính mô-đun của số phức $z-w$
\choice
{$\sqrt{3}$}
{\True $\sqrt{2}$}
{$\sqrt{5}$}
{$2$}
\loigiai{
Số phức $z$ có điểm biểu diễn trong mặt phẳng $Oxy$ là $M\left( 2\,;\,1 \right)$ $\Rightarrow z=2+i$\\
Số phức $w$ có điểm biểu diễn trong mặt phẳng $Oxy$ là $N\left( 1\,;\,2 \right)$ $\Rightarrow w=1+2i$\\
Khi đó: $\left| z-w \right|=\left| \left( 2+i \right)-\left( 1+2i \right) \right|=\left| 1-i \right|=\sqrt{1^2+{{\left( -1 \right)}^2}}=\sqrt{2}$
.}
\end{ex}

\begin{ex}%[2D4Y1-2]%
Số phức $z$ thỏa mãn $\bar{z}=1-2i$ được biểu diễn trên mặt phẳng tọa độ bởi điểm nào sau đây?
\choice
{ $Q\left( -1;-2 \right)$}
{\True $M\left( 1;2 \right)$}
{ $P\left( -1;2 \right)$}
{ $N\left( 1;-2 \right)$}
\loigiai{
Vì $\bar{z}=1-2i\Rightarrow z=1+2i$.\\
Do đó điểm biểu diễn số phức $z$ là $\left( 1;2 \right)$.}
\end{ex}

\begin{ex}%[2D4Y1-2]
Trên mặt phẳng tọa độ $Oxy$, điểm biểu diễn số phức $z=(1+2i)^2$ là điểm nào dưới đây?
\choice
{\True $P(-3;4)$}
{$Q(5;4)$}
{$N(4;-3)$}
{$M(5;4)$}
\loigiai{
Ta có $z=(1+2i)^2=1+4i+4i^2=-3+4i\Rightarrow$ điểm biểu diễn số phức $ z=(1+2i)^2$ là điểm $P(-3;4)$.
}
\end{ex}

\begin{ex}%[2D4B1-1]%
Cho số phức $z$ thỏa mãn $z \cdot\bar{z}=16$ giá trị của $|z|$ tương ứng bằng
\choice
{$16$}
{\True $4$}
{$\pm4$}
{$2$}
\loigiai{
Ta có $z \cdot\bar{z}=16\Leftrightarrow|z|^2=16\Leftrightarrow z=4$.
}
\end{ex}

\begin{ex}%[Đề thi thử - Trường THPT chuyên Lương Thế Vinh - Đồng Nai - Lần 1 - 2018]%[2D4B1-2]%[Kim Minh Bui - 12EX8]%
\immini{Điểm $M$ trong hình vẽ bên là điểm biểu diễn của số phức $z$. Tìm phần thực và phần ảo của số phức $z$.
\choice
{Phần thực là $4$ và phần ảo là $3$}
{Phần thực là $3$ và phần ảo là $4i$}
{\True Phần thực là $3$ và phần ảo là $4$}
{Phần thực là $4$ và phần ảo là $3i$}}
{\begin{tikzpicture}
\draw[->] (-0.5,0) -- (2.5,0) node[below] {$x$};
\draw[->] (0,-0.5) -- (0,2.5) node[right] {$y$};
\node at (0,0) [below left] {\footnotesize $O$};
\node at (1.5,0) [below] {$3$};
\node at (0,2) [left] {$4$};
\node at (1.5,2) [above right] {$M$};
\draw[dashed]
(0,2)--(1.5,2)--(1.5,0);
\end{tikzpicture}}
\loigiai{Điểm $ M(3;4) $ là điểm biểu diễn của số phức $ z=3+4i $.}
\end{ex}

\begin{ex}%[2D4Y2-1]%
cho hai số phức $ z=4+3i$ và $ w=1-i$. Số phức $ z-w$ bằng
\choice
{ $5+2i$}
{ $7-i$}
{\True $3+4i$}
{ $-3-4i$}
\loigiai{
$ z-w=4+3i-( 1-i )=3+4i$.
}
\end{ex}

\begin{ex} %[2D4Y2-2]%
Cho số phức $z$ thỏa mãn $z=3-4 i$. Điểm biểu diễn số phức $u=i z$ trên mặt phẳng phức có tọa độ tương ứng là
\choice
{\True $(4 ; 3)$}
{$(3 ;-4)$}
{$(3 ; 4)$}
{$(3 ;-3)$}
\loigiai{
Ta có $u=i z=3i+4$, nên điểm biểu diễn số phức $u=i z$ trên mặt phẳng phức có tọa độ tương ứng là $(4;3)$.
}
\end{ex}

\begin{ex}%[TT Chuyên Quốc Học Huế L2, Huế 2018]%[Trần Nhân Kiệt, dự án EX10]%[2D4B2-1]%
Cho hai số phức $z$ và $z'$. Trong các mệnh đề sau, mệnh đề nào {\bf sai}?
\choice
{\True $|z+z'|=| z|+| z'|$}
{$| z\cdot z'|=| z|\cdot| z'|$}
{$\overline{z}\cdot\overline{z'}=\overline{z\cdot z'}$}
{$\overline{z}+\overline{z'}=\overline{z+z'}$}
\loigiai{
Mệnh đề $| z+z'|=| z|+|z'|$ {\bf sai} vì với $z=1+i$ và $z'=1-i$ thì
\begin{align*}
&|z+z'|=|(1+i)+(1-i)|=|2|=2\\
&|z|+|z'|=|1+i|+|1-i|=2\sqrt{2}\\
\Rightarrow & |z+z'|\neq |z|+|z'|.
\end{align*}
}
\end{ex}

\begin{ex}%[2D4Y3-1]%
Cho hai số phức $z_1=1+3i$ và $z_2=2i-3$. Số phức $\dfrac{z_1}{z_2}$ bằng
\choice
{$\dfrac{-3-11i}{13}$}
{$-\dfrac{3-11i}{13}$}
{$\dfrac{3+11i}{13}$}
{\True $\dfrac{3-11i}{13}$}
\loigiai{
Ta có $\dfrac{z_1}{z_2}=\dfrac{1+3i}{2i-3}=\dfrac{3-11i}{13}$.
}
\end{ex}

\begin{ex}%[2D4B3-2]%[TheHung Nguyen]%
Cho số phức $z$ thỏa mãn $(1-2i)z=(4+3i)(2-z)$. Giá trị $|z|$ bằng
\choice
{$\dfrac 5{\sqrt{13}}$}
{\True $\dfrac{5\sqrt{26}}{13}$}
{$2\sqrt 3$}
{$\dfrac{5\sqrt 5}2$}
\loigiai{
Ta có $(1-2i)z=(4+3i)(2-z)\Leftrightarrow z((1-2i)+(4+3i))=2(4+3i)$.\\
Suy ra $z=\dfrac{2(4+3i)}{(1-2i)+(4+3i)}$.\\
Vậy $|z|=\left|\dfrac{2(4+3i)}{(1-2i)+(4+3i)}\right|=\dfrac{5\sqrt{26}}{13}$.
}
\end{ex}

\begin{ex}%[2D4B3-2]%
Số phức $z$ thỏa mãn điều kiện $z=4+3\overline{z}$ có mô-đun bằng
\choice
{\True $\sqrt{5}$}
{$2$}
{$-2$}
{$4$}
\loigiai{
Đặt $z=a+bi$, $(a;b\in\mathbb{R})$.\\
Ta có $z=4+3\overline{z} \Leftrightarrow \heva{&a=4+3a\\&b=4-3b} \Leftrightarrow \heva{&a=-2\\&b=1.}$\\
Vậy $|z|=\sqrt{(-2)^2+1^2}=\sqrt{5}$.
}
\end{ex}

\begin{ex}%[Vũ Ngọc Phát]%[TDM]%[2D4B3-4]%
Cho số phức $z$ thỏa mãn $|z-1+2i|=|z+3i-2|$. Trên mặt phẳng tọa độ $Oxy$ tập hợp điểm biểu diễn số phức $z$ là
\choice
{\True đường thẳng $x-y-4=0$}
{đường thẳng $x+y-8=0$}
{đường thẳng $2x-y+3=0$}
{đường thẳng $x^2+(y-2)^2=36$}
\loigiai{
Gọi $M(x;y)$, $A(1;-2)$, $B(2;-3)$ lần lượt là điểm biểu diễn các số phức $z$, $z_1=1-2i$, $z_2=2-3i$.\\
Theo đề ta có $|z-z_1|=|z-z_2|$. Suy ra $MA=MB$. Do đó tập hợp các điểm $M$ là đường trung trực $d$ của đoạn thẳng $AB$.\\
Đường thẳng $d$ đi qua trung điểm $I\left(\dfrac{3}{2};-\dfrac{5}{2}\right)$ của $AB$ và nhận $\overrightarrow{AB}=(1;-1)$ làm véc-tơ pháp tuyến nên có phương trình là
$$\left(x-\dfrac32\right) - \left(y+\dfrac52\right)=0 \Leftrightarrow x-y-4=0.$$
Vậy tập hợp điểm biểu diễn số phức $z$ thỏa yêu cầu bài toán là đường thẳng $x-y-4=0$.
}
\end{ex}

\begin{ex}%[Số phức - Tư duy mở]%[Hồng Trường Sơn]%[2D4B4-1]%
Cho hai số phức $z_1;z_2$ là nghiệm của phương trình $z^2-2z+3=0$. Gọi $A$ và $B$ là hai điểm biểu diễn số phức $z_1$ và $z_2$. Khoảng cách $AB$ bằng
\choice
{$2\sqrt3$}
{$4\sqrt2$}
{\True $2\sqrt2$}
{$4\sqrt5$}
\loigiai{
Ta có $z^2-2z+3=0 \Leftrightarrow \hoac{&z_1=1+i\sqrt2 \Rightarrow A\left(1;\sqrt2\right)\\ &z_2=1-i\sqrt2 \Rightarrow B\left(1;-\sqrt2\right).}$\\
Vậy $AB=2\sqrt2$.
}
\end{ex}

\begin{ex}%[Số phức - Tư duy mở]%[Hồng Trường Sơn]%[2D4Y4-1]%
Có bao nhiêu số phức $z$ thỏa mãn phương trình $z^2-4z+3=0$?
\choice
{$1$}
{$3$}
{\True $2$}
{$0$}
\loigiai{
Ta có $z^2-4z+3=0 \Leftrightarrow \hoac{&z=1\\ &z=3.}$\\
Vậy phương trình đã cho có $2$ nghiệm phức.
}
\end{ex}

\begin{ex}%[Nguyễn Thành Tiến- Dự án Tổng ôn LVĐ 2022]%[2H3Y1-1]%
Trong KG $Oxyz$, cho $\vec{OA}=3\vec{i}+4\vec{j}-5\vec{k}$. Tọa độ điểm $A$ là
\choice
{\True $A\left(3;4;-5\right)$}
{$A\left(-3;4;5\right)$}
{$A\left(3;4;5\right)$}
{$A\left(-3;-4;5\right)$}
\loigiai{
Ta có $A\left(3;4;-5\right)$.
}
\end{ex}

\begin{ex}%[2H3B1-3]%
Trong không gian Oxyz, cho mặt cầu $(S): x^{2}+y^{2}+z^{2}-2 x-4 y-6 z+5=0$. Diện tích của mặt cầu $(S)$ bằng
\choice
{$42 \pi$}
{\True $36 \pi$}
{$9 \pi$}
{$12 \pi$}
\loigiai{
Mặt cầu đã cho có tâm $I(1;2;3)$, bán kính $R=\sqrt{1^2+2^2+3^2-5}=3$.\\
Diện tích mặt cầu là $4 \pi 3^2=36 \pi$
}
\end{ex}

\begin{ex}%[2H3Y2-2]%
Trong không gian $O x y z$, mặt phẳng $(P): 2 x-6 y-8 z+1=0$ có một véc-tơ pháp tuyến là
\choice
{$(-1 ;-3 ; 4)$}
{$(1 ; 3 ; 4)$}
{\True $(1 ;-3 ;-4)$}
{$(1 ;-3 ; 4)$}
\loigiai{
$(P): 2 x-6 y-8 z+1=0 \Rightarrow \text{ VTPT } \vec{n}=(2;-6;-8)=2(1;-3;-4)$ nên $\text{ VTPT } \vec{n}=(1;-3;-4)$.
}
\end{ex}

\begin{ex}%[2H3Y2-4]%
Trong KG $Oxyz$, cho mặt phẳng $(P)\colon 2 x+y-z-6=0$. Điểm nào sau đây không thuộc mặt phẳng $(P)$ ?
\choice
{$(0 ; 3 ;-3)$}
{$(3 ; 0 ; 0)$}
{$(3 ; 1 ; 1)$}
{\True $(3 ; 2 ;-2)$}
\loigiai{
Vì $2\cdot 3+2-(-2)-6=4 \neq 0$ nên điểm có tọa độ $(3 ; 2 ;-2)$ không thuộc mặt phẳng $(P)$.}
\end{ex}

\begin{ex}%[2H3B2-3]%
Trong KG $Oxyz$, mặt phẳng $(P)$ đi qua hai điểm $A(1;-1;2)$, $B(2;1;1)$ và vuông góc với mặt phẳng $(Q)\colon x+y+z+1=0$ có phương trình là
\choice
{\True $3x-2y-z-3=0$}
{$x+y+z-2=0$}
{$x-y=0$}
{$3x-2y-z+3=0$}
\loigiai{
Ta có $\overrightarrow{AB}=(1;2;-1)$.\\
Mặt phẳng $(Q)$ có véc-tơ pháp tuyến là $\overrightarrow{n}_Q=(1;1;1)$.\\
Mặt phẳng $(P)$ đi qua hai điểm $A$, $B$ và vuông góc với mặt phẳng $(Q)$ có véc-tơ pháp tuyến
\[\overrightarrow{n}=\left[\overrightarrow{AB};\overrightarrow{n}_Q\right]=(3;-2;-1).\]
Phương trình mặt phẳng $(P)$ là
\[3\cdot (x-1)-2\cdot (y+1)-1\cdot (z-2)=0\Leftrightarrow 3x-2y-z-3=0.\]
}
\end{ex}

\begin{ex}%[2H3B2-3]%
Trong KG $Oxyz$, cho mặt phẳng $(\alpha)$ đi qua $A(1;-2;3)$ và song song mặt phẳng $(Oxy)$ thì phương trình mặt phẳng $(\alpha)$ là
\choice
{$x-1=0$}
{$x+2y+z=0$}
{$y+2=0$}
{\True $z-3=0$}
\loigiai{
Vì $(\alpha)\parallel(Oxy):z=0$ nên $(\alpha): z+d=0$ với $d\not=0$.\\
Mặt phẳng $(\alpha)$ đi qua $A(1;-2;3)$ nên $A(1;-2;3)\in (\alpha):z+d=0\Rightarrow3+d=0\Rightarrow d=-3$ (nhận).\\
Vậy $(\alpha):z-3=0$.
}
\end{ex}

\begin{ex}%[2H3Y3-1]
Trong KG $Oxyz$ , một véctơ chỉ phương của đường thẳng
$\Delta\colon \heva{& x=1\\ & y=2+2t\\&     z=1-3t}$ là
\choice
{$\vec{u}=\left(0;2;3\right)$}
{$\vec{u}=\left(1;2;-3\right)$}
{\True $\vec{u}=\left(0;2;-3\right)$}
{$\vec{u}=\left(1;2;1\right)$}
\loigiai{
Đường thẳng $\Delta\colon \heva{& x=1\\ & y=2+2t\\&     z=1-3t}$ có một véctơ chỉ phương là $\vec{u}=\left(0;2;-3\right)$.}
\end{ex}

\begin{ex}%[2H3Y3-6]%
Trong KG $Oxyz$, mặt phẳng nào trong các mặt phẳng sau song song với đường thẳng
$(d)\colon  \dfrac{x+1}{1}=\dfrac{y+2}{2}=\dfrac{z+3}{3}$.
\choice
{$(P)\colon  z+2 y+3 z=5$}
{$(Q)\colon  3 x-y-2 z=5$}
{\True $(\alpha)\colon  3 x-3 y+z=5$}
{$(\beta)\colon  3 x-3 y+z=0$}
\loigiai{
Đường thẳng $(d)$ có một vectơ chỉ phương là vectơ $\vec{u}=(1 ; 2 ; 3)$, và đi qua điểm $A(-1 ;-2 ;-3)$.\\
Mặt phẳng $(\alpha)$ và $(\beta)$ đều có vectơ pháp tuyến vuông góc với vectơ $\vec{u}$, và $A \notin(\alpha)$, nên mặt phẳng $(\alpha)$ song song đường thẳng $(d)$.
}
\end{ex}

\begin{ex}%[TexBook12-Quyen3, Biên soạn thầy Nguyễn Ngọc Nguyên - Phản biện thầy Nguyễn Sĩ Đạt]% %[2H3Y3-3]%
Trong KG $Oxyz$, cho đường thẳng $d: \heva{&x=1+2t \\ &y=2-t \\ &z=-2+2t}$ ($t\in\mathbb{R}$) và điểm $M(1;2;m)$. Tìm giá trị của tham số $m$ để điểm $M$ thuộc đường thẳng $d$.
\choice
{$m=2$}
{\True $m=-2$}
{$m=1$}
{$m=0$}
\loigiai{
$$ M\in d \Leftrightarrow \heva{&1=1+2t \\ &2=2-t \\ &m=-2+2t} \Leftrightarrow \heva{&t=0 \\ &t=0 \\ &m=-2.}$$
Vậy $m=-2$ thì điểm $M$ thuộc đường thẳng $d$.
}
\end{ex}

\begin{ex}%[2H3B3-2]%
Cho hàm số $f(x)$ có đạo hàm $f'(x)=x(x-1)(x+4)^3$, $\forall x\in\mathbb{R}$. Số điểm cực tiểu của hàm số đã cho là
\choice
{\True $2$}
{$3$}
{$4$}
{$1$}
\loigiai{
Ta có $f'(x)=0\Leftrightarrow x(x-1)(x+4)^3=0\Leftrightarrow \left[\begin{array}{l}
x=0\\
x=1\\
x=-4
\end{array}\right.$\\
Bảng biến thiên:
\begin{center}
\begin{tikzpicture}
\tkzTabInit[nocadre,lgt=1.5,espcl=3]{$x$/0.8,$f'(x)$/0.8,$f(x)$/2.5}{$-\infty$,$-4$,$0$,$1$,$+\infty$}
\tkzTabLine{,-,0,+,0,-,0,+}
\tkzTabVar{+/$+\infty$,-/,+/,-/,+/$+\infty$}
\end{tikzpicture}
\end{center}
Dựa vào bảng biến thiên ta thấy hàm số đã cho có 2 điểm cực tiểu.
}
\end{ex}

\begin{ex}%[Hình học Oxyz, Tư duy mở]%[Lê Hồng Phi]%[2H3B3-6]%
Trong không gian tọa độ $Oxyz$, cho hai đường thẳng $\Delta_1\colon \dfrac{x-3}{2}=\dfrac{y+2}{3}=\dfrac{z-1}{1}$ và đường thẳng $\Delta_2\colon \dfrac{x+2}{1}=\dfrac{y-a}{m}=\dfrac{z-b}{n}$. Nếu hai đường thẳng $\Delta_1$, $\Delta_2$ trùng nhau thì ta có $(a+b+m+n)$ tương ứng bằng
\choice
{$11$}
{$7$}
{$-12$}
{\True $-9$}
\loigiai{
Để $\Delta_1\equiv\Delta_2$ thì trước hết phải có $\dfrac{1}{2}=\dfrac{m}{3}=\dfrac{n}{1}\Leftrightarrow\heva{& m=\dfrac{3}{2} \\ & n=\dfrac{1}{2}.}$\\
Tiếp đến, điểm $M(3;-2;1)\in\Delta_1$ thì $M\in\Delta_2\Rightarrow\dfrac{3+2}{1}=\dfrac{-2-a}{\frac{3}{2}}=\dfrac{1-b}{\frac{1}{2}}\Rightarrow\heva{& a=-\dfrac{19}{2} \\ & b=-\dfrac{3}{2}.}$\\
Vậy $a+b+m+n=-9$.
}
\end{ex}

\begin{ex}%[2D3K2-3]%
Giá trị của tích phân $I=\displaystyle \int\limits_{\frac{\pi}{6}}^{\frac{\pi}{3}} \cos x \ln \left(\sin x\right) \mathrm{d}x$ tương ứng bằng
\choice
{$\dfrac{\sqrt{3}}{2}-1$}
{\True $\dfrac{\sqrt{3}\ln 3 +\left(2-2\sqrt{3}\right)\left(\ln 2+1\right)}{4}$}
{$\dfrac{3\sqrt{3}-2}{4}$}
{$\sqrt{3}\ln 3 -2\ln 2+1$}
\loigiai{
Đây là tích phân từng phần. Đặt $\heva{
&u=\ln \left(\sin x\right) \\ & \mathrm{d}v =\cos x \mathrm{d}x} \Rightarrow \heva{
&\mathrm{d}u=\cos x \cdot \dfrac{\mathrm{d}x}{x}\\
&v = \sin x.}$\\
Suy ra $I=\displaystyle \int\limits_{\frac{\pi}{6}}^{\frac{\pi}{3}}\cos x \ln \left(\sin x\right)\mathrm{d}x=\sin x \ln \left(\sin x\right) \bigg|_{\frac{\pi}{6}}^{\frac{\pi}{3}} - \int\limits_{\frac{\pi}{6}}^{\frac{\pi}{3}}\sin x \cdot \cos x \dfrac{\mathrm{d}x}{\sin x}.$\\
Hay $I=\dfrac{\sqrt{3}}{2}\ln \dfrac{\sqrt{3}}{2}-\dfrac{1}{2}\ln \dfrac{1}{2} \displaystyle \int\limits_{\frac{\pi}{6}}^{\frac{\pi}{3}}\cos x \mathrm{d}x$.\\
Dẫn đến $I=\dfrac{\sqrt{3}}{2}\ln \dfrac{\sqrt{3}}{2}-\dfrac{1}{2}\ln \dfrac{1}{2}-\sin x \bigg|_{\frac{\pi}{6}}^{\frac{\pi}{3}}=\dfrac{\sqrt{3}\ln 3+\left(2-2\sqrt{3}\right)\left(\ln 2+1\right)}{4}.$
}
\end{ex}

\begin{ex}%[2D3K2-2]%
Tích phân $\displaystyle \int\limits_0^{2020}x\left(x-1\right)\left(x-2\right)...\left(x-2020\right)\mathrm{d}x$ bằng
\choice
{$2020!$}
{\True $0$}
{$2019!$}
{$\dfrac{2020!}{2}$}
\loigiai{
Đặt $t=2020-x \Rightarrow \mathrm{d}t = -\mathrm{d}x.$ Đổi cận $\heva{
&x=0 \Rightarrow t=2020 \\&x=2020 \Rightarrow t=0.}$\\
Suy ra $I=\displaystyle \int\limits_0^{2020} x\left(x-1\right)\left(x-2\right)...\left(x-2020\right)\mathrm{d}x = \displaystyle \int\limits_{2020}^{0}\left(2020-t\right)\left(2019-t\right)\left(2018-t\right)...\left(-t\right)\left(-\mathrm{d}t\right)$.\\
Dẫn đến $I=-\displaystyle \int\limits_0^{2020}t\left(t-1\right)\left(t-2\right)...\left(t-2020\right)\mathrm{d}t=-I \Leftrightarrow I=0$.
}
\end{ex}

\begin{ex}%[2D3G2-4]%
Cho hàm số $f(x)$ có đạo hàm cấp hai xác định trên $\mathbb{R}$ thỏa mãn hệ thức $$f''(x)+f(x)=2f'(x)+x \text{ } \forall x \in \mathbb{R}.$$ Biết $f'(0)=1$, $f(0)=2$. Giá trị $\displaystyle \int\limits_{0}^{1}f(x) \mathrm{\,d}x$ bằng
\choice
{$3$}
{$\dfrac{7}{2}$}
{\True $\dfrac{5}{2}$}
{$\dfrac{4}{3}$}
\loigiai{\noindent
$f''(x)+f(x)=2f'(x)+x \Leftrightarrow \left(f''(x)-f'(x)\right)-\left(f'(x)-f(x)\right)=x$.\\
Đặt $u(x)=f'(x)-f(x) \Rightarrow u'(x)=f''(x)-f'(x)$. Suy ra
\begin{eqnarray*}
&&\left(f''(x)-f'(x)\right)-\left(f'(x)-f(x)\right)=x \\
&\Leftrightarrow &u'(x)-u(x)=x \\
&\Leftrightarrow &\mathrm{e}^{-x}\left(u'(x)-u(x)\right)=\mathrm{e}^{-x}x\\
&\Leftrightarrow &\left(\mathrm{e}^{-x}\cdot u(x)\right)'=\mathrm{e}^{-x}x\\ &\Leftrightarrow &\mathrm{e}^{-x} \cdot u(x)=\displaystyle \int \mathrm{e}^{-x}x \mathrm{\,d}x=-\mathrm{e}^{-x}(x+1)+C.
\end{eqnarray*}
Thay $x=0 \Rightarrow u(0)=f'(0)-f(0)=1-2=-1$.\\
Suy ra $\mathrm{e}^{-0}\cdot u(0)=-\mathrm{e}^{-0}(0+1)+C \Rightarrow C=0$.
\begin{eqnarray*}
&&\mathrm{e}^{-x}\cdot u(x)=-\mathrm{e}^{-x}(x+1)\\
&\Leftrightarrow & u(x)=-x-1=f'(x)-f(x) \\
&\Leftrightarrow &\mathrm{e}^{-x} \left(f'(x)-f(x)\right)=\mathrm{e}^{-x}(-x-1)\\
&\Leftrightarrow &\left[\mathrm{e}^{-x}f(x)\right]'=\mathrm{e}^{-x}(-x-1)\\
&\Leftrightarrow &\mathrm{e}^{-x}f(x)=\displaystyle \int \mathrm{e}^{-x}(-x-1)\mathrm{\,d}x=\mathrm{e}^{-x}(x+2)+C_1.
\end{eqnarray*}
Thay $x=0$, $f(0)=2 \Rightarrow \mathrm{e}^{-0}f(0)=\mathrm{e}^{-0}\cdot (0+2)+C_1 \Rightarrow C_1=0$.\\
Suy ra $\mathrm{e}^{-x}f(x)=\mathrm{e}^{-x}(x+2) \Leftrightarrow f(x)=x+2$.\\
Do đó $\displaystyle \int\limits_{0}^{1}f(x) \mathrm{\,d}x=\displaystyle \int\limits_{0}^{1} (x+2) \mathrm{\,d}x=\dfrac{5}{2}$.

}
\end{ex}

\begin{ex}%[Trần Nhân Kiệt]%[Dự án TexBook12]%[2D4G5-1]%
Cho hai số phức $z_1$, $z_2$ khác $0$ thỏa mãn $\dfrac{z_1}{z_2}$ là số thuần ảo và $\left| z_1-z_2 \right|=10$. Giá trị lớn nhất của $\left| z_1 \right|+\left| z_2 \right|$ bằng
\choice
{$10$}
{\True $10\sqrt{2}$}
{$10\sqrt{3}$}
{$20$}
\loigiai{
Cách 1:\\
Vì $\dfrac{z_1}{z_2}$ là số thuần ảo nên $\dfrac{z_1}{z_2}=ai\Leftrightarrow z_1=aiz_2$.\\
Ta có $\left| z_1-z_2 \right|=10\Leftrightarrow \left| aiz_2-z_2 \right|=10\Leftrightarrow\left| z_2 \right|\cdot\left| ai-1 \right|=10\Leftrightarrow \left| z_2 \right|\sqrt{1+a^2}=10\Leftrightarrow \left| z_2 \right|=\dfrac{10}{\sqrt{1+a^2}}$.\\
Từ $z_1=aiz_2$
$\left| z_1 \right|=\left| aiz_2 \right|=\dfrac{10\left| a \right|}{\sqrt{1+a^2}}$.\\
Do đó $\left| z_1 \right|+\left| z_2 \right|=\dfrac{10\left| a \right|}{\sqrt{1+a^2}}+\dfrac{10}{\sqrt{1+a^2}}$ $=\dfrac{10\left( 1+\left| a \right| \right)}{\sqrt{1+a^2}}$ $\le \dfrac{10\sqrt{\left( 1+1 \right)\left( 1+a^2 \right)}}{\sqrt{1+a^2}}$ $\le 10\sqrt{2}$.\\
Đẳng thức xảy ra $ a=\pm 1$ $z_1=\pm iz_2$.\\
Vậy $ \max\left( \left| z_1 \right|+\left| z_2 \right| \right)=10\sqrt{2}$.\\
Cách 2:\\
Đặt $z_1=a_1+b_1i$, $z_2=a_2+b_2i$.\\
Gọi $A$, $B$ lần lượt là điểm biểu diễn của $z_1$, $z_2$\\
$\Rightarrow A\left( a_1\,;\,b_1 \right)$, $B\left( a_2\,;\,b_2 \right)$\\
$\Rightarrow \overrightarrow{OA}\left( a_1\,;\,b_1 \right)\,,\,\overrightarrow{OB}\left( a_2\,;\,b_2 \right)$.\\
Ta có $\dfrac{z_1}{z_2}=\dfrac{z_1\cdot \overline{z_2}}{\left| z_2 \right|^2}=\dfrac{\left( a_1+b_1i \right)\left( a_2-b_2i \right)}{\left| z_2 \right|^2}$ là số thuần ảo $\Leftrightarrow a_1a_2+b_1b_2=0\Leftrightarrow \overrightarrow{OA}\cdot\overrightarrow{OB}=0$.\\
Suy ra $\Delta OAB$ vuông tại $O$.\\
Khi đó $\left| z_1-z_2 \right|=AB=10$.\\
Ta có $\left| z_1 \right|+\left| z_2 \right|=OA+OB\le \sqrt{\left( OA^2+OB^2 \right)\cdot\left( 1^2+1^2 \right)}=\sqrt{2AB^2}=10\sqrt{2}$.\\
Đẳng thức xảy ra $OA=OB$.\\
Vậy $ \max\left( \left| z_1 \right|+\left| z_2 \right| \right)=10\sqrt{2}$.
}

\end{ex}

\begin{ex}%[Phạm Hoàng Điệp]%[2H3K2-3]%
Trong KG $Oxyz$, cho điểm $H(a;b;c)$ với $a$, $b$, $c>0$. Mặt phẳng $(P)$ chứa điểm $H$ và lần lượt cắt các trục $Ox$, $Oy$, $Oz$ tại $A,B,C$ thỏa mãn $H$ là trực tâm của tam giác $ABC$. Phương trình của mặt phẳng $(P)$ là
\choice 
{$\dfrac{x}{a^{2}}+\dfrac{y}{b^{2}}+\dfrac{z}{c^{2}}=\dfrac{ab+bc+ca}{abc}$}
{$\dfrac{x}{a}+\dfrac{y}{b}+\dfrac{z}{c}=3$}
{\True $ax+by+cz-a^{2}-b^{2}-c^{2}=0$}
{$a^{2}x+b^{2}y+c^{2}z-a^{3}-b^{3}-c^{3}=0$}
\loigiai{
\textbf{Cách 1}\\
Gọi $A\left(x_{0};0;0\right)$, $B\left(0 ; y_{0} ; 0\right)$, $C\left(0 ; 0 ; z_{0}\right)$.\\
Khi đó mặt phẳng $(P)$ có phương trình theo đoạn chắn là: $\dfrac{x}{x_{0}}+\dfrac{y}{y_{0}}+\dfrac{z}{z_{0}}=1$.\\
Ta có $\overrightarrow{AH}=\left(a-x_{0} ; b ; c\right)$, $\overrightarrow{BC}=\left(0 ;-y_{0} ; z_{0}\right)$, $\overrightarrow{BH}=\left(a ; b-y_{0} ; c\right)$, $\overrightarrow{AC}=\left(-x_{0} ; 0 ; z_{0}\right)$.\\
Vì $H$ là trực tâm tam giác $A B C$ nên ta có hệ\\
$ \heva{&\overrightarrow{AH} \cdot \overrightarrow{BC}=0 \\ &\overrightarrow{BH} \cdot \overrightarrow{AC}=0 \\ &H \in(ABC)&} \Leftrightarrow \heva{&-by_{0}+cz_{0}=0 \\ &-ax_{0}+cz_{0}=0 \\&\dfrac{a}{x_{0}}+\dfrac{b}{y_{0}}+\dfrac{c}{z_{0}}=1} \Leftrightarrow \heva{&y_{0}=\dfrac{c}{b}z_{0} \\ &x_{0}=\dfrac{c}{a}z_{0} \\ &\dfrac{a}{\dfrac{c}{a}z_{0}}+\dfrac{b}{\dfrac{b}{c}z_{0}}+\dfrac{c}{z_{0}}=1}\Leftrightarrow \heva{&y_{0}=\dfrac{a^{2}+b^{2}+c^{2}}{b} \\& x_{0}=\dfrac{a^{2}+b^{2}+c^{2}}{a} \\& z_{0}=\dfrac{a^{2}+b^{2}+c^{2}}{c}.}$\\
Thay vào phương trình mặt phẳng $(P)$ ta được $\dfrac{ax}{a^{2}+b^{2}+c^{2}}+\dfrac{by}{a^{2}+b^{2}+c^{2}}+\dfrac{cz}{a^{2}+b^{2}+c^{2}}=1$.\\
Hay $(P)\colon ax+by+cz-a^{2}-b^{2}-c^{2}=0.$\\
\textbf{Cách 2}\\
Ta chứng minh được $OH \perp(ABC)$ hay $OH \perp(P)$.\\
Do đó mặt phẳng $(P)$ qua $H$ và nhận $\overrightarrow{OH}(a ; b ; c)$ làm véc-tơ pháp tuyến có phương trình là \begin{center}
$a(x-a)+b(y-b)+c(z-c)=0 \Leftrightarrow ax+by+cz-a^{2}-b^{2}-c^{2}=0 .$
\end{center}
}
\end{ex}

\begin{ex}%[2H3K3-7]%
Trong không gian tọa độ $Oxyz$, cho hai điểm $A(3;1;1)$, $B(7;3;9)$ và mặt phẳng $(P)\colon x+y+z+3=0$. Điểm $M(x;y;x)\in (P)$ sao cho $\left|\vv{MA}+\vv{MB}\right|$ đạt giá trị nhỏ nhất. Giá trị $x+y+z$ bằng
\choice
{\True $-3$}
{$3$}
{$0$}
{$2$}
\loigiai{
Gọi $I$ là điểm thỏa mãn $\vv{IA}+\vv{IB}=\vv{0}$ thì $I$ là trung điểm của $AB$. Suy ra $I(5;2;5)$.\\
Khi đó, $\left|\vv{MA}+\vv{MB}\right|=\left|\vv{MI}+\vv{IA}+\vv{MI}+\vv{IB}\right|=2\left|\vv{MI}\right|$.\\
$\left|\vv{MA}+\vv{MB}\right|$ đạt giá trị nhỏ nhất $\Leftrightarrow\left|\vv{MI}\right|$ đạt giá trị nhỏ nhất $\Leftrightarrow M$ là hình chiếu vuông góc của $I$ lên mặt phẳng $(P)$.\\
Gọi $d$ là đường thẳng đi qua $I$ và vuông góc với $(P)$. Suy ra $M=d\cap (P)$.\\
$d\colon\heva{&x=5+t\\&y=2+t\\&z=5+t}\Rightarrow M(5+t;2+t;5+t)$.\\
$M\in (P)$ nên $5+t+2+t+5+t+3=0\Leftrightarrow t=-5\Rightarrow M(0;-3;0)$.\\
Vậy $x+y+x=-3$.\\
{\bf Trắc nghiệm} Vì $M(x;y;z)\in (P)$ nên $x+y+z+3=0\Leftrightarrow x+y+z=-3$.
}
\end{ex}


\Closesolutionfile{ans}